\documentclass[11pt]{article}
\usepackage[margin=2cm]{geometry}
\usepackage[T1]{fontenc}
\usepackage{helvet}
\usepackage{hyperref}
\usepackage{listings}
\usepackage{longtable}
\usepackage{booktabs}

\renewcommand{\familydefault}{\sfdefault}

\lstset{
  basicstyle=\ttfamily\small,
  breaklines=true,
  columns=fullflexible
}

\title{Hydrogenic-demo Documentation}
\author{Hydrogenic-demo contributors}
\date{September 2026}

\begin{document}
\maketitle

\section{Overview}

Hydrogenic-demo provides a compact C++ example for hydrogenic atom calculations, together with a Python interface for interactive work and notebook-based demonstrations. The C++ code builds the \texttt{hydrogenic\_atom\_demo} executable and exercises the routines in \texttt{Hydrogenic.vX} together with numerical utilities from \texttt{Tools}.

\section{Project layout}

\begin{longtable}{p{0.34\linewidth}p{0.58\linewidth}}
\toprule
Path & Purpose \\
\midrule
\texttt{main.cpp} & Example program that constructs hydrogenic atoms and prints representative atomic data. \\
\texttt{Hydrogenic.vX/Atom.*} & Core \texttt{Atom}, \texttt{Atomic\_Shell}, \texttt{Electron\_Level}, and \texttt{Gas\_of\_Atoms} classes. \\
\texttt{Hydrogenic.vX/Oscillator\_strength.*} & Bound-bound oscillator strength and transition support. \\
\texttt{Hydrogenic.vX/Photoionization\_cross\_section.*} & Photoionization cross sections and Gaunt factors. \\
\texttt{Hydrogenic.vX/Quadrupole\_lines\_HI.*} & Electric quadrupole line rates for hydrogenic atoms. \\
\texttt{Hydrogenic.vX/Rec\_Phot\_BB.*} and \texttt{Recombination\_Integrals.*} & Recombination and photon interaction helpers. \\
\texttt{Tools/} & Shared constants, integration routines, line profiles, and simple numerical helpers. \\
\texttt{python/} & Installable Python package and Jupyter notebook examples. \\
\texttt{Documentation/} & PDF documentation and editable TeX source. \\
\texttt{Makefile} and \texttt{Makefile.in} & Build entry points for the C++ demo executable. \\
\bottomrule
\end{longtable}

\section{C++ example}

The demo constructs hydrogenic atoms with a selected number of complete shells. Each shell contains angular momentum levels with \texttt{l = 0 ... n - 1}. The code can report level degeneracies, ionization energies, ionization frequencies, total decay rates, lifetimes, selected bound-bound transition data, optional electric quadrupole transitions, photoionization cross sections, Gaunt factors, and simple rescaling behavior for the fine-structure constant and electron mass.

The current \texttt{main.cpp} includes hydrogen with \texttt{Z = 1}, 5 shells, quadrupole lines enabled, and Storey-Hummer recombination interpolation enabled. It also includes hydrogenic oxygen, O VIII, with \texttt{Z = 8}, 30 shells, quadrupole lines enabled, and Storey-Hummer recombination interpolation enabled.

\section{Scientific background}

The hydrogenic atom setup was developed for detailed calculations of cosmological recombination. This requires multi-shell hydrogenic atoms with resolved angular momentum sub-states, dipole transitions, quadrupole transitions, photoionization cross sections, and recombination coefficients. The C++ example in \texttt{main.cpp} cites the main background papers for this use case.

The physics model is the standard one-electron Coulomb problem for a nucleus of charge \texttt{Z}, with reduced-mass corrections included in the Rydberg scaling. Radiative quantities are evaluated perturbatively in the Born approximation: unperturbed nonrelativistic Coulomb bound and continuum states are coupled by the leading electromagnetic multipole operators. This makes the routines appropriate for hydrogenic ions such as H I and O VIII, but not for detailed multi-electron structure.

The bound-bound dipole rates and oscillator strengths use the Storey \& Hummer recurrence relations for reduced radial matrix elements. The implementation reorganizes these recursions for numerical stability at large \texttt{n}, using root-rescaled starting values and delayed products so that very small radial integrals are not repeatedly multiplied into underflow or cancellation. The bound-free photoionization cross sections and Gaunt factors use the corresponding Storey \& Hummer bound-free recurrences; the Kramers reference cross section follows Karzas \& Latter. The electric quadrupole routines combine the same hydrogenic radial machinery with quadrupole matrix-element relations from Hey and the recombination application discussed by Grin \& Hirata.

\section{Dependencies}

The C++ build expects a C++17 compiler with OpenMP support, GSL headers and libraries, Boost headers, and GNU \texttt{make}. The Python package additionally uses Python 3.9 or newer, \texttt{pybind11}, and \texttt{numpy}; the notebook examples use \texttt{jupyter} and \texttt{matplotlib}.

On macOS with Homebrew:

\begin{lstlisting}[language=bash]
brew install gcc gsl boost
\end{lstlisting}

\section{Build and run}

\begin{lstlisting}[language=bash]
make
make help
make h
make usage
make list
make targets
make run
make clean
make update-modules
\end{lstlisting}

\texttt{make --help} and \texttt{make -h} are handled by GNU make itself; use \texttt{make help} for project-specific help.

The executable is written to \texttt{./hydrogenic\_atom\_demo}; intermediate C++ build products are written under \texttt{./build/}. Python extension build products are written under \texttt{python/build/} and \texttt{python/src/hydrogenic/}. \texttt{make clean} removes both the C++ and Python build products.

\section{Python package and Jupyter}

Use \texttt{python3}, not \texttt{python}, from the repository root. Since this repository has a \texttt{python/} folder, some shell setups can otherwise resolve \texttt{python} incorrectly and report \texttt{Permission denied}.

To build the Python extension and run the notebook with existing Homebrew/JupyterLab installations, you need \texttt{numpy}, \texttt{matplotlib}, \texttt{jupyterlab} or \texttt{notebook}, build-time \texttt{setuptools} and \texttt{pybind11}, and the compiled dependencies GCC with OpenMP, GSL, and Boost headers.

On macOS with Homebrew:

\begin{lstlisting}[language=bash]
brew install gcc gsl boost python-setuptools pybind11 numpy matplotlib jupyterlab
\end{lstlisting}

Build the local Python extension:

\begin{lstlisting}[language=bash]
make python-build
make python
make py
\end{lstlisting}

Open the demo notebook:

\begin{lstlisting}[language=bash]
make run-jupyter
\end{lstlisting}

\texttt{make run-jupyter} runs \texttt{make python-build} first, then opens the notebook with an available Jupyter command (\texttt{jupyter lab}, \texttt{jupyter-lab}, or \texttt{jupyter-notebook}) and sets \texttt{PYTHONPATH} to include \texttt{python/src}. This lets the notebook import the local \texttt{hydrogenic} extension without installing it into JupyterLab's private Python environment.

The Makefile workflow avoids installing the package into your system Python. If you deliberately want a pip install into a Homebrew-managed Python, Homebrew may require:

\begin{lstlisting}[language=bash]
python3 -m pip install --user --break-system-packages ./python
\end{lstlisting}

If needed, override compiler and dependency paths explicitly:

\begin{lstlisting}[language=bash]
CXX=/opt/homebrew/bin/g++-16 \
GSL_INC_PATH=/opt/homebrew/include \
GSL_LIB_PATH=/opt/homebrew/lib \
BOOST_INC_PATH=/opt/homebrew/include \
python3 -m pip install ./python
\end{lstlisting}

Basic Python use:

\begin{lstlisting}[language=Python]
import numpy as np
from hydrogenic import hydrogen, oxygen

h = hydrogen(shells=5)
print(h.level(2, 1))
print(h.transition(2, 1, 1, 0))

o = oxygen(shells=30)
nu0 = o.level(1, 0)["ionization_frequency_hz"]
nu = np.geomspace(1.001, 10.0, 100)*nu0
data = o.photoionization(1, 0, nu)
print(data["sigma_cm2"])
\end{lstlisting}

Photoionization cross sections require \texttt{recombination\_mode=1} or \texttt{recombination\_mode=2}; this is the default used by the Python convenience constructors.

\section{Python notebooks}

The notebook examples live in \texttt{python/notebooks/}. They illustrate the same routines from Python, including level data, line transitions, and photoionization cross-section plots for hydrogen and hydrogenic oxygen.

\section{References}

\begin{enumerate}
\item J. Chluba and R. M. Thomas, ``Towards a complete treatment of the cosmological recombination problem'', MNRAS 412, 748, 2011. \url{https://ui.adsabs.harvard.edu/abs/2011MNRAS.412..748C/abstract}
\item J. Chluba, J. A. Rubino-Martin, and R. A. Sunyaev, ``Cosmological hydrogen recombination: populations of the high level sub-states'', MNRAS 374, 1310, 2007. \url{https://ui.adsabs.harvard.edu/abs/2007MNRAS.374.1310C/abstract}
\item J. Chluba and Y. Ali-Haimoud, ``CosmoSpec: Fast and detailed computation of the cosmological recombination radiation from hydrogen and helium'', MNRAS 456, 3494, 2016. \url{https://ui.adsabs.harvard.edu/abs/2016MNRAS.456.3494C/abstract}
\item P. J. Storey and D. G. Hummer, ``Fast computer evaluation of radiative properties of hydrogenic systems'', Computer Physics Communications 66, 129, 1991. \url{https://ui.adsabs.harvard.edu/abs/1991CoPhC..66..129S/abstract}
\item W. J. Karzas and R. Latter, ``Electron Radiative Transitions in a Coulomb Field'', ApJS 6, 167, 1961. \url{https://ui.adsabs.harvard.edu/abs/1961ApJS....6..167K/abstract}
\item J. D. Hey, ``On the determination of radial matrix elements for high-n transitions in hydrogenic atoms and ions'', J. Phys. B: At. Mol. Opt. Phys. 39, 2641, 2006. \url{https://doi.org/10.1088/0953-4075/39/12/003}
\item D. Grin and C. M. Hirata, ``Cosmological hydrogen recombination: The effect of extremely high-n states'', Phys. Rev. D 81, 083005, 2010. \url{https://arxiv.org/abs/0911.1359}
\end{enumerate}

\end{document}
